\section{Kesimpulan}
Rekonstruksi tanpa FL Studio menghasilkan lima komponen dari 2219 not
yang penempatannya mengikuti Playlist. Hasil menggambarkan komponen
aransemen dengan SoundFont pengganti, bukan reproduksi identik audio
asli. Bass paling terkonsentrasi pada frekuensi rendah; tuba mempunyai
timing serupa dengan pitch satu oktaf lebih tinggi; piano dan gitar
memuat lapisan not yang berubah; brass berupa dua kelompok akor
dengan proporsi harmonik tinggi yang besar.

FFT menunjukkan frekuensi berkontribusi global kuat dan distribusi
energi. Sebagai contoh, 97,61\% energi bass berada di bawah 250 Hz,
sedangkan 55,70\% energi brass berada di atas 1000 Hz. Angka ini
mendukung perbedaan komposisi frekuensi, bukan penetapan keras suara
perseptual atau fundamental otomatis. RMS seluruh timeline juga
perlu dibedakan dari RMS aktif, terutama pada brass yang banyak diam.

STFT menambahkan lokasi waktu: pergantian not, perubahan lapisan
melodi dan akor yang bertahan dapat dilihat langsung. Jendela 1024
lebih terlokalisasi dalam waktu, sedangkan 4096 memisahkan garis
frekuensi lebih jelas, seperti pada brass sekitar 1,4--1,6 kHz.
Kompromi ini sesuai Heisenberg; nilai 0,513117 berasal dari model
Hann kontinu, bukan pengukuran ketelitian pendeteksi audio.
Parseval mendukung konsistensi energi dengan galat relatif di bawah
$7\times10^{-15}$.

Analisis kanal audio menambahkan perbedaan pola aktivitas: DRUMS
aktif sepanjang sebagian besar rekaman, sedangkan VOCALS mempunyai
kejadian kuat pada beberapa interval dengan total aktivitas sekitar
14,2 s menurut ambang yang digunakan. DRUMS lebih didominasi
frekuensi rendah, sementara VOCALS terutama memuat energi menengah
dan tinggi. WAV keduanya menyimpan sampel hasil dekode MP3 secara
identik, sehingga konversi tidak meningkatkan kualitas sumber maupun
ketelitian FFT/STFT; hasil DRUMS yang dihitung ulang tetap konsisten
dengan analisis sebelumnya.

Perubahan RMS, spektrum lokal dan distribusi pita mendukung deskripsi
nonstasioner, tetapi bukan uji statistik formal. Penelitian belum
menentukan kesamaan timbre terhadap FLEX/DirectWave, karakter umum
semua instrumen sejenis, atau sinkronisasi audio sumber penuh.
DRUMS dan VOCALS tetap menjadi pembanding independen dan tidak
dimasukkan ke campuran lima instrumen.

\begin{thebibliography}{99}
\raggedright
\bibitem{kuliah4}\textit{Kuliah 4: Transformasi Fourier}. Bahan kuliah
Analisis Fourier dan Wavelet, tanpa tahun. Folder \texttt{Materi},
\texttt{Kuliah-04-AFW.pdf}, Bagian 4.2 dan 4.5, hlm.~52--53, 58--59.
\bibitem{kuliah5}\textit{Kuliah 5: Sampling dan Analisis Fourier Diskret}.
Bahan kuliah Analisis Fourier dan Wavelet, tanpa tahun.
\texttt{KuliahAFW-05.pdf}, Bagian 5.1 dan 5.3--5.5, hlm.~65, 69--75.
\bibitem{kuliah6}\textit{Kuliah 6: Analisis Waktu-Frekuensi}. Bahan kuliah
Analisis Fourier dan Wavelet, tanpa tahun. \texttt{KuliahAFW-06.pdf},
Bagian 6.3--6.7, hlm.~89--100; tugas Bagian 6.10, hlm.~104.
\bibitem{youtube}Rick Ruyz. \textit{[ FREE FLP ] Bad Piggies - Theme REMAKE
FL STUDIO STOCK PLUGINS}. YouTube, tanpa tahun.
\url{https://www.youtube.com/watch?v=DdIcPpjQXV4}.
Judul dan kanal diperiksa 22 Oktober 2026; hubungan dengan paket lokal
mengikuti keterangan penyedia data.
\bibitem{flp}Image-Line. \textit{FL Studio Project File (*.flp)}.
\url{https://www.image-line.com/fl-studio-learning/fl-studio-online-manual/html/fformats_open_flp.htm}.
Diakses 22 Oktober 2026.
\bibitem{numpy}NumPy Developers. \textit{Discrete Fourier Transform}.
\url{https://numpy.org/doc/stable/reference/routines.fft.html}.
Diakses 22 Oktober 2026.
\bibitem{hann}SciPy Developers. \textit{scipy.signal.windows.hann}.
\url{https://docs.scipy.org/doc/scipy/reference/generated/scipy.signal.windows.hann.html}.
Diakses 22 Oktober 2026.
\bibitem{pyflp}PyFLP. \textit{Documentation: Patterns, Notes, Arrangements}.
\url{https://pyflp.readthedocs.io/}. Versi parser yang digunakan 2.2.1.
\bibitem{playlist}PyFLP. \textit{Missing playlist data in FL Studio 2025
project files}, Issue 200. \url{https://github.com/demberto/PyFLP/issues/200}.
Bandingkan dokumentasi pemeriksaan format independen:
\url{https://github.com/dawhubapp/flpdiff/blob/main/docs/fl-format/flp-format-spec.md}.
Diakses 22 Oktober 2026.
\bibitem{fluidsynth}FluidSynth Developers. \textit{FluidSynth 2.6.1}.
\url{https://github.com/FluidSynth/fluidsynth/releases/tag/v2.6.1}.
Mesin sintesis SoundFont; diakses 22 Oktober 2026.
\bibitem{soundfont}S.~Christian Collins. \textit{GeneralUser GS 2.0.3}.
\url{https://www.schristiancollins.com/generaluser.php}.
Sumber \url{https://github.com/mrbumpy409/GeneralUser-GS}; lisensi v2.0
disertakan bersama dataset. Diakses 22 Oktober 2026.
\bibitem{soundfile}Bastian Bechtold dan kontributor.
\textit{python-soundfile 0.13.1: Reading and Writing Sound Files}.
\url{https://python-soundfile.readthedocs.io/en/latest/}.
Dokumentasi pembacaan sampel dan penyimpanan WAV float;
diakses 22 Oktober 2026.
\end{thebibliography}

\appendix
\section{Normalisasi STFT dan Model Hann}\label{app:normalisasi}
Rincian ini melengkapi uraian utama. Dengan $H$ sampel pergeseran,
DFT potongan berjendela ialah
\[
 Z_{q,m}=\sum_{r=0}^{M-1}f_{qH+r}w_r e^{-2\pi i mr/M},\quad
 \tau_q=\frac{qH+M/2}{f_s},\quad \omega_m=\frac{2\pi mf_s}{M}.
\]
Hubungan dengan STFT materi:
\[
 V_wf(\tau_q,\omega_m)\approx\frac{\Delta t}{\sqrt{2\pi}}
 e^{-i\omega_mqH\Delta t}Z_{q,m}.
\]
Faktor fase tidak mengubah spektrogram. Kerapatan daya satu sisi
yang digunakan untuk grafik adalah
\[
 P_{q,m}=\frac{c_m|Z_{q,m}|^2}{f_s\sum_rw_r^2},\quad
 \Spec_wf\approx\frac{\Delta t^2}{2\pi}|Z_{q,m}|^2
 =\frac{\sum_rw_r^2}{2\pi f_sc_m}P_{q,m}.
\]
Grafik merupakan penskalaan spektrogram per energi jendela dan per
Hz, bukan nilai mentah spektrogram materi. Nilai dB adalah
$10\log_{10}(P/P_0)$ dengan $P_0=1$ satuan amplitudo kuadrat per Hz,
bukan dB SPL. Pemeriksaan daya bingkai memakai
\[
 \sum_mP_{q,m}\frac{f_s}{M}
 =\frac{\sum_r|f_{qH+r}w_r|^2}{\sum_rw_r^2}.
\]
Penjumlahan semua piksel spektrogram secara langsung bukan energi
rekaman, karena tumpang tindih, bobot jendela dan kisi perlu
diperhitungkan.

Definisi penyebaran Heisenberg menggunakan pusat energi
$t_w=\int t|w|^2dt/\|w\|_2^2$ dan
$\omega_w=\int\omega|\widehat w|^2d\omega/\|\widehat w\|_2^2$, kemudian
\[
 \sigma_t^2=\frac{\int(t-t_w)^2|w(t)|^2dt}{\|w\|_2^2},\qquad
 \sigma_\omega^2=\frac{\int(\omega-\omega_w)^2|\widehat w(\omega)|^2d\omega}{\|\widehat w\|_2^2}.
\]
Hann kontinu adalah $w_a(t)=\frac12(1+\cos(2\pi t/a))$ untuk
$|t|\leq a/2$, dan nol di luarnya. Kedua pusat nol karena simetri.
Integrasi memberi
\[
 \|w_a\|_2^2=\frac{3a}{8},\quad
 \sigma_t=a\sqrt{\frac1{12}-\frac5{8\pi^2}},\quad
 \|w_a'\|_2^2=\frac{\pi^2}{2a}.
\]
Sifat turunan Fourier dan Plancherel kemudian memberi
\[
 \sigma_\omega=\frac{2\pi}{\sqrt3a},\qquad
 \sigma_t\sigma_\omega=\sqrt{\frac{\pi^2}{9}-\frac56}
 \approx0{,}513117.
\]
Dengan $a=M/f_s$, $\sigma_t$ sekitar 3,284 dan 13,138 ms;
$\sigma_\omega$ sekitar 156,228 dan 39,057 rad/s. Nilai ini
penyebaran jendela kontinu, bukan durasi dukungan, hop, atau galat
waktu kejadian audio.

\section{Reproduksi dan Pemeriksaan Hasil}
Notebook \texttt{analisis\_instrumen.ipynb} memuat seluruh kode
persiapan perangkat, ekstraksi FLP, rekonstruksi MIDI dan WAV,
analisis FFT dan STFT, serta validasi di dalam sel-selnya.
Eksekusi berurutan menghasilkan CSV not dan Playlist,
grafik, tabel dan pemutar cuplikan. Hasil baru berada di
\texttt{hasil/instrumen}, grafik di \texttt{gambar/instrumen},
dan audio pendek di \texttt{hasil/audio/cuplikan}. Sampling WAV adalah
44.100 Hz. Data sumber ditempatkan di \texttt{data/sumber}, sedangkan
asal paket dan video dicatat pada \texttt{data/sumber\_data.json}.

Bagian analisis kanal sumber dalam notebook yang sama mengonversi
DRUMS dan VOCALS ke \texttt{hasil/audio/kanal}, memeriksa kesamaan
sampel hasil dekode, dan menyimpan hasil baru di
\texttt{hasil/kanal\_audio} serta \texttt{gambar/kanal\_audio}.
RMS, energi dan puncak DRUMS dari WAV dibandingkan kembali dengan
perhitungan langsung pada dekode MP3, sehingga pemeriksaan tidak
bergantung pada arsip hasil terdahulu. Kesamaan sampel juga diperiksa
untuk VOCALS. CSV not, Playlist dan audit sampel berada pada
\texttt{hasil/rekonstruksi}; MIDI berada pada \texttt{hasil/midi}.

Untuk mereproduksi hasil, buka notebook dengan kernel Python 3.11,
pasang dependensi dari \texttt{requirements.txt} bila diperlukan,
kemudian jalankan semua sel secara berurutan. Bagian awal notebook
menjelaskan susunan berkas dan parameter eksekusi. FluidSynth portabel
dan dependensi lokal berada dalam \texttt{perangkat}; SoundFont dan
lisensinya berada dalam \texttt{data/soundfont}. Pengunduhan diperlukan
hanya jika perangkat tersebut belum tersedia. Rekonstruksi dapat
diulang dengan mengaktifkan \texttt{REBUILD\_AUDIO}; analisis instrumen
dan kanal sumber dihitung ulang oleh sel masing-masing.

Paper ini memuat penjelasan ilmiah, interpretasi grafik, audit
rekonstruksi dan batas kesimpulan. Notebook memuat kode dan petunjuk
komputasinya. Seluruh gambar dikelompokkan dalam \texttt{gambar}
dan seluruh keluaran numerik serta audio dalam \texttt{hasil}.
PDF dapat dikompilasi dari \texttt{ETS AFW.tex} dengan menjalankan
\texttt{pdflatex} dua kali setelah gambar dan tabel tersedia.
